\newpage

% ######################################################################
% PART XIII: EXTRA Q&A + PROJECT REFERENCE
% ######################################################################
\part{Extra Q\&A + Project Reference}

% ######################################################################
% SECTION 53: 20 MORE INTERVIEW QUESTIONS
% ######################################################################
\section{20 More Interview Questions --- Mixed Topics}

\subsection{AI/ML Concepts}

\begin{interviewbox}
\textbf{Q76: ``Overfitting kya hai? LLM mein hota hai?''}

\textbf{Answer:} Overfitting: model training data memorize karta hai, generalize nahi. LLM fine-tuning mein ho sakta hai (data memorize). Is project mein overfitting relevant nahi (hum fine-tune nahi karte) --- par \textbf{prompt overfitting} hota hai: prompts specific examples par tune ho jate hain aur real-world fail. Isliye benchmark on unseen topics zaroori.
\end{interviewbox}

\begin{interviewbox}
\textbf{Q77: ``Precision vs recall? Is project mein kaunsa metric?''}

\textbf{Answer:} Precision: jo bole wo kitna sahi (false positives kam). Recall: jo sahi tha wo kitna pakda (false negatives kam). RAG retrieval: precision = retrieved chunks relevant hain? recall = saare relevant chunks aaye? Is project ke benchmark mein depth/verifiability quality metrics hain --- retrieval metrics (recall@k) production mein add karne ki plan hai.
\end{interviewbox}

\begin{interviewbox}
\textbf{Q78: ``Gradient descent intuition?''}

\textbf{Answer:} Loss function ko minimize karte hue weights update: $w = w - \alpha \nabla L$. Learning rate $\alpha$: bada = fast par overshoot; chhota = slow par stable. Is project mein hum trained model use karte hain --- gradient descent training ka core hai, inference forward pass only.
\end{interviewbox}

\begin{interviewbox}
\textbf{Q79: ``Batch vs online learning?''}

\textbf{Answer:} Batch: pura dataset ek saath train. Online/incremental: data aate hi update. Is project mein RAG ek tarah ka ``online knowledge'' hai --- naye PDFs embed hokar turant available, bina model retrain kiye. Ye hi RAG vs fine-tuning ka core difference hai.
\end{interviewbox}

\begin{interviewbox}
\textbf{Q80: ``Confusion matrix?''}

\textbf{Answer:} TP/FP/TN/FN ka table. Is project mein: fact-checker ke VERIFIED/UNVERIFIED decisions ko confusion matrix mein rakha ja sakta hai --- VERIFIED par source support nahi = false positive (hallucination slipped). Ye evaluation maturity dikhata hai.
\end{interviewbox}

\subsection{Product \& Design Thinking}

\begin{interviewbox}
\textbf{Q81: ``Is system ka target user kaun hai?''}

\textbf{Answer:} Teen user types: (1) \textbf{Knowledge worker} --- deep research reports chahiye (Streamlit); (2) \textbf{AI assistant user} --- Claude/Cursor se research tool call (MCP); (3) \textbf{Developer} --- programmatic access + benchmark (Django API). Har interface apne user ke consumption pattern ke liye designed hai.
\end{interviewbox}

\begin{interviewbox}
\textbf{Q82: ``Is system ki moat kya hai?''}

\textbf{Answer:} (Honest) Core moat nahi hai --- LLMs + Tavily + ChromaDB sab commodity hai. Moat: (1) \textbf{Workflow design} --- fact-check + revision loop ka specific orchestration; (2) \textbf{Evaluation data} --- benchmark results jo accumulate hote hain; (3) \textbf{UX} --- 3 interfaces ka integration. Ye bolo to product thinking dikhti hai.
\end{interviewbox}

\begin{interviewbox}
\textbf{Q83: ``Ek feature drop karna ho to kaunsa?''}

\textbf{Answer:} (Prioritization) MCP server drop kar sakte hain --- niche use case (AI assistants ko tools chahiye), audience chhota. Core (pipeline + UI + benchmark) rakho. Ya PDF upload drop karo agar koi use nahi karta --- data se decide karo (analytics dekho), assumptions se nahi.
\end{interviewbox}

\begin{interviewbox}
\textbf{Q84: ``User feedback kaise collect karoge?''}

\textbf{Answer:} (1) \textbf{Quantitative} --- benchmark scores, latency, usage analytics; (2) \textbf{Qualitative} --- report quality rating (thumbs up/down), open feedback; (3) \textbf{Interviews} --- 5 users se 30-min sessions. Is project mein benchmark hi quantitative feedback loop hai.
\end{interviewbox}

\begin{interviewbox}
\textbf{Q85: ``Is project ko open source karna chahiye?''}

\textbf{Answer:} Trade-offs: \textbf{Pro:} credibility, community contributions, portfolio visibility. \textbf{Con:} maintenance burden, secrets leak risk (env gitignored par careful), competition. Personal project ke liye: haan, open source + good README --- ye interview portfolio ke liye strong signal hai.
\end{interviewbox}

\subsection{Miscellaneous Technical}

\begin{interviewbox}
\textbf{Q86: ``JSON vs XML vs YAML?''}

\textbf{Answer:} JSON: lightweight, universal, no comments --- API standard. XML: verbose, schema validation (XSD) --- legacy/enterprise. YAML: human-friendly, comments, nested --- config files. Is project mein: API responses JSON (DRF default), config .env (key-value), MCP schemas JSON, aur ye document LaTeX.
\end{interviewbox}

\begin{interviewbox}
\textbf{Q87: ``REST vs GraphQL?''}

\textbf{Answer:} REST: resource-based endpoints, simple, cacheable --- is project mein. GraphQL: client-controlled queries, over-fetching/under-fetching solve --- par complexity + caching hard. Is scale par REST right choice; GraphQL tab jab client data needs vary karein (mobile + web + dashboard).
\end{interviewbox}

\begin{interviewbox}
\textbf{Q88: ``Synchronous vs asynchronous programming?''}

\textbf{Answer:} Sync: ek kaam khatam, phir agla --- simple, blocking. Async: I/O wait mein doosra kaam --- efficient. Is project mein pipeline sync hai (176s block) --- production mein async (task queue) because users wait nahi karte. Python async: asyncio/await + event loop; concurrency (not parallelism) for I/O.
\end{interviewbox}

\begin{interviewbox}
\textbf{Q89: ``Docker kya hai? Container vs VM?''}

\textbf{Answer:} Container: app + dependencies packaged, OS kernel shared --- light, fast start. VM: full OS per instance --- heavy, isolated. Is project ke liye Docker: \texttt{python:3.11-slim} base, requirements install, streamlit/django CMD --- reproducible environments (``works on my machine'' khatam).
\end{interviewbox}

\begin{interviewbox}
\textbf{Q90: ``Git branching strategy?''}

\textbf{Answer:} Simple: main (stable) + feature branches + PR review + merge. Production: main protected, release branches, tags. Is project mein main par direct work hua (solo) --- team mein feature branches + PRs. Interview mein trunk-based vs git-flow dono batao.
\end{interviewbox}

\begin{interviewbox}
\textbf{Q91: ``Logging vs print?''}

\textbf{Answer:} Logging: levels (DEBUG/INFO/WARNING/ERROR), configurable, timestamps, structured. Print: ad-hoc, no control. Is project mein agents \texttt{print} karte hain (sabse basic) --- production mein \texttt{logging} module + structured JSON logs + log aggregation. Interview line: ``print demo ke liye, logging production ke liye.''
\end{interviewbox}

\begin{interviewbox}
\textbf{Q92: ``Environment variables vs config files?''}

\textbf{Answer:} Env vars: secrets + environment-specific (keys, URLs) --- gitignored, 12-factor. Config files: non-secret defaults (chunk size, model name) --- committed, versioned. Is project mein: secrets env/secrets; defaults config/setting.py mein with env override. Ye separation security + clarity dono deta hai.
\end{interviewbox}

\begin{interviewbox}
\textbf{Q93: ``API key leak ho jaye to kya karoge?''}

\textbf{Answer:} (Incident response) (1) \textbf{Revoke immediately} --- provider dashboard se key delete; (2) Naya key generate + rotate; (3) \textbf{Source hunt} --- git history se remove (git filter-repo), secrets scanning; (4) Damage assessment --- usage logs check karo unauthorized calls; (5) Prevention --- .env gitignore, secret managers, pre-commit hooks. Ye structured response impress karta hai.
\end{interviewbox}

\begin{interviewbox}
\textbf{Q94: ``Documentation kyu important?''}

\textbf{Answer:} (1) Onboarding --- naya developer 10 min mein setup (README + run.md); (2) \textbf{Knowledge preservation} --- 6 mahine baad khud bhi bhool jaoge; (3) \textbf{Quality signal} --- interviewers README dekhte hain; (4) \textbf{API docs} --- consumers ke liye contract. Is project mein: README (setup + API table), run.md (preview ops), ye 150-page document (concepts + code + interview).
\end{interviewbox}

\begin{interviewbox}
\textbf{Q95: ``Koi feature ship karte waqt testing strategy?''}

\textbf{Answer:} (1) \textbf{Unit tests} --- logic (evaluator heuristics); (2) \textbf{Integration} --- API + mocked pipeline; (3) \textbf{E2E} --- real run (manual, paid); (4) \textbf{Benchmark} --- quality regression; (5) \textbf{Monitoring} --- production metrics post-ship. Is project mein 16 tests + benchmark + real-run verification --- multi-layer strategy.
\end{interviewbox}

% ######################################################################
% SECTION 54: DEPENDENCY REFERENCE
% ######################################################################
\section{Dependency Reference --- Har Package Kya Karta Hai}

\begin{longtable}{p{4.2cm} p{5.4cm} p{4.8cm}}
\toprule
\textbf{Package} & \textbf{Role} & \textbf{Project mein kahan} \\
\midrule
langchain-core & Messages, prompts & agents/*.py \\
langchain-groq & ChatGroq LLM client & agents ka get\_llm \\
langchain-community & PyPDFLoader, Chroma & rag/vector\_store.py \\
langchain-text-splitters & RecursiveCharacterTextSplitter & rag ingestion \\
langchain-huggingface & HuggingFaceEmbeddings & rag embeddings \\
langgraph & StateGraph orchestration & graph/workflow.py \\
tavily-python & Web search client & researcher \\
sentence-transformers & MiniLM model & embeddings (transitive) \\
chromadb & Vector store & rag/vector\_store.py \\
streamlit & Web UI & app.py, pages/ \\
mcp / fastmcp & MCP server & mcp\_server.py \\
django & Web framework & backend/ \\
djangorestframework & REST API layer & backend/api/ \\
django-cors-headers & CORS policy & backend settings \\
python-dotenv & .env loading & config/setting.py \\
requests & HTTP client & pages/1\_Benchmark.py \\
pandas & DataFrames + charts & pages/1\_Benchmark.py \\
pytest & Testing & backend/api/tests/ \\
pytest-django & Django test integration & pytest.ini \\
\bottomrule
\end{longtable}

\begin{memorybox}
\textbf{Interview line:} ``Dependencies ka har package ka role pata hai --- langchain ecosystem orchestration/retrieval ke liye, Django API ke liye, Streamlit UI ke liye, pytest quality ke liye. Koi dependency bina kaam ke nahi hai.''
\end{memorybox}

% ######################################################################
% SECTION 55: PROJECT METRICS
% ######################################################################
\section{Project Metrics --- Numbers Jo Yaad Rakho}

\begin{longtable}{p{6cm} p{3cm} p{5cm}}
\toprule
\textbf{Metric} & \textbf{Value} & \textbf{Context} \\
\midrule
Agents & 4 & researcher, analyst, fact-checker, writer \\
Interfaces & 3 & Streamlit, MCP, Django REST \\
Tests & 16 & auth(6) + research(4) + benchmark(6) \\
Test time & \textasciitilde9s & Fast suite = daily use \\
Pipeline time & \textasciitilde176s & End-to-end real run \\
Sources per run & \textasciitilde15 & 3 sub-queries x 5 results \\
Fact-check passes & 1st pass & Solid-state run, revision 0 \\
Max revisions & 3 & Loop guard \\
Chunk size / overlap & 1000 / 200 & Chars \\
Embedding dims & 384 & all-MiniLM-L6-v2 \\
Model & GPT-OSS-120B & Via Groq \\
Endpoints & 6 & health, register, login, execute, benchmark, history \\
Benchmark dims & 2 & Depth + verifiability (0-10) \\
Benchmark verdicts & 3 & single/multi superior + tie \\
DB tables & 2 + auth & ResearchSession, BenchmarkEvaluation \\
\bottomrule
\end{longtable}

\begin{storybox}
\textbf{Final message:}

Ye document ab poori tarah taiyaar hai. 150+ pages, har concept, har file, har logic --- Hinglish mein, story format mein. \textbf{Ab aapki baari hai:} ise padho, samjho, aur sabse zaroori --- \textbf{khud bol ke practice karo}. Interview success = understanding + confidence + communication. Teeno is document se ready hain.

\textbf{Shubhkamnayein!} \faHandPeace
\end{storybox}
